\documentclass[11pt,a4paper]{article}
\usepackage[T1]{fontenc}
\usepackage{lmodern,amsmath,amssymb,amsthm,mathtools,booktabs,microtype}
\usepackage[margin=27mm]{geometry}
\usepackage[colorlinks=true,linkcolor=blue,citecolor=blue,urlcolor=blue]{hyperref}
\setlength{\emergencystretch}{2em}
\newtheorem{theorem}{Theorem}[section]
\newtheorem{lemma}[theorem]{Lemma}
\newtheorem{proposition}[theorem]{Proposition}
\newcommand{\Q}{\mathbb Q}
\newcommand{\Z}{\mathbb Z}
\DeclareMathOperator{\Nm}{N}
\title{A two-coordinate pole test of coefficient height four}
\author{}\date{17 November 2026}
\begin{document}\maketitle

\begin{abstract}
A polynomial of total degree nine and coefficient height four gives a
pole test with two unrestricted rational parameters and one quadratic
root. Its rational solutions exclude negative odd valuations of the
input at every prime. For positive inputs below one, with square
reduced denominator and a specified $2$-adic square class, Burungale's
simultaneous-twist theorem supplies rational witnesses. Their projection
to the parameter plane is Zariski dense, which permits simultaneous
avoidance of the denominator and radicand zeros at every character of a
finite point algebra. The pole block retains two free coordinates; no
smaller defining-variable bound follows from the block alone.
\end{abstract}

\section{The equation and its local implication}

Set
\begin{equation}\label{eq:radicand}
 P_d(v,s)=d(d+v^2)^2s^4-d(d-v^2)^2(1-d)^2.
\end{equation}
For a rational input $d$, use rational parameters $v,s$ and impose
\begin{equation}\label{eq:pole}
 Y^2=P_d(v,s).
\end{equation}
No nonzero condition is needed for the local implication below.
The equation is monic quadratic in $Y$. Its radicand has total degree
nine, twelve monomials, and coefficient height four.
For example, the prepared values
\[
 d=\frac{112}{841},\qquad v=\frac4{29},\qquad s=\frac{27}{29},\qquad
 Y=\frac{653184}{20511149}
\]
satisfy the equation with nonzero ordinate. They arise from the
point $(28,98)$ on $E_{21}$ and $1-d=(27/29)^2$.

\begin{theorem}\label{thm:local}
If \eqref{eq:pole} has a rational solution, then $v_p(d)$ is not a
negative odd integer at any prime $p$. In particular its reduced
denominator is a square.
\end{theorem}
\begin{proof}
Suppose $n=v_p(d)<0$ is odd and use the convention $v_p(0)=+\infty$.
Since $n$ and $2v_p(v)$ have different parity whenever both are finite,
\[
 v_p(d+v^2)=v_p(d-v^2)=r:=\min(n,2v_p(v)).
\]
Also $v_p(1-d)=n$. The two summands in \eqref{eq:radicand} have
valuations
\[
 n+2r+4v_p(s),\qquad 3n+2r.
\]
If $s\ne0$, both are odd and they are unequal: equality would imply
$2v_p(s)=n$. If $s=0$, only the second summand remains, and it has
odd valuation. In either case $P_d(v,s)$ is nonzero with odd valuation,
contradicting its being a rational square. This argument applies at
every prime, including two, and permits $v=0$ and $Y=0$.
\end{proof}

\section{The arithmetic input}

For $l\ne0$ put $E_l:y^2=x^3-l^2x$.
We use the following part of Burungale's Theorem B~\cite{Burungale}.
Let $\beta,\eta$ be coprime positive squarefree integers such that
\[
 \beta\equiv7\pmod8,\qquad \eta\equiv1\pmod8.
\]
Suppose $\eta$ is a square in $\Q_p$ at each prime dividing $\beta$,
and $\beta$ is a square in $\Q_p$ at each prime dividing $\eta$.
Then there is a positive squarefree integer $n\equiv5\pmod8$ for
which $E_n$ and $E_{n\eta}$ have rank one and rational points with
abscissa in $\beta\Q^{\times2}$. The source proves additional
conclusions that are not needed here. The application below verifies
every hypothesis of this theorem; it does not independently reprove
the source's arithmetic proof.

\begin{theorem}\label{thm:density}
Assume the arithmetic input above. Suppose $d\in\Q$ has square
reduced denominator and
\begin{equation}\label{eq:prepared}
 0<d<1,\qquad d\in7\Q_2^{\times2},\qquad v_2(d)\ge4.
\end{equation}
The projection to the $(v,s)$ plane of the rational solutions of
\eqref{eq:pole} is Zariski dense, including over $\overline\Q$.
\end{theorem}
\begin{proof}
Write $d=A/D^2$ in lowest terms with $A,D>0$. Let $\beta$ be the
positive squarefree representative of $d$ and let $\eta$ be that of
$1-d=(D^2-A)/D^2$. The two integer numerators are coprime, so
$\beta$ and $\eta$ are coprime. At each odd prime dividing $\beta$,
$1-d$ is a unit with residue one; hence $\eta$ is a local square.
At each odd prime dividing $\eta$, the value $d$ is a unit with
residue one; hence $\beta$ is a local square.

The $2$-adic square class of $d$ gives $\beta\equiv7\pmod8$.
Its valuation is even and at least four, so $1-d\equiv1\pmod{16}$
and $\eta\equiv1\pmod8$. Thus the arithmetic input applies, also
when $\eta=1$. Fix the resulting $n$ and write $1-d=\eta c^2$,
$c\in\Q^\times$. Scaling the point on $E_{n\eta}$ by
$(x,y)\mapsto(c^2x,c^3y)$ gives a point on $E_{n(1-d)}$ with the
same abscissa square class.

Each of $E_n$ and $E_{n(1-d)}$ has infinitely many rational points
in the abscissa square class $d$. Indeed the abscissa square-class
homomorphism is unchanged upon adding an element of $2E_l(\Q)$,
and each curve has positive rank. Exclude the finitely many torsion
points and write two independently chosen points as
\[
 R=(da^2,aU)\in E_n(\Q),\qquad
 R'=(db^2,bV)\in E_{n(1-d)}(\Q),
\]
where $abUV\ne0$. Their positive abscissas exceed the corresponding
positive curve parameters. In particular, define
\[
 e=\frac{n}{da^2}\in(0,1),\qquad
 p_0=\frac{U}{da^2},\qquad
 s=\frac ba,\qquad q_0=\frac{V}{da^2},\qquad
 v=\frac{p_0}{1+e}.
\]
The curve equations give
\[
 p_0^2=d(1-e^2),\qquad
 q_0^2=d\bigl(s^4-e^2(1-d)^2\bigr).
\]
The first relation yields
\[
 e=\frac{d-v^2}{d+v^2},\qquad p_0=\frac{2dv}{d+v^2}.
\]
It follows that $Y=(d+v^2)q_0$ gives a solution of \eqref{eq:pole}.
Moreover $vsY(d-v^2)(d+v^2)\ne0$ by the choices just made.

Varying $R$ gives infinitely many values of $a^2$, hence of $e$ and
$v^2=d(1-e)/(1+e)$, and therefore infinitely many $v$. For each
fixed $R$, varying $R'$ gives infinitely many $s$. A polynomial that
vanishes on all witness pairs therefore vanishes as a polynomial in
$s$ for infinitely many $v$; all its coefficients vanish identically.
This proves the asserted density.
\end{proof}

\section{Guards in a finite algebra}

\begin{proposition}\label{prop:guards}
Let $\mathcal A$ be a finite etale $\Q$-algebra with
$d\in\mathcal A^\times$. Suppose a rational character of
$\mathcal A$ sends $d$ to an input satisfying \eqref{eq:prepared}.
Actual rational witnesses may be chosen so that all factors
\[
 s,\quad d-v^2,\quad d+v^2,\quad P_d(v,s)
\]
are nonzero at every geometric character of $\mathcal A$.
Any additional finite collection of guards with nonzero polynomial
norms in $v,s$ may be avoided at the same time.
\end{proposition}
\begin{proof}
At a character taking $d$ to $c\ne0$, the two quadratic factors
are nonzero polynomials in $v$, and $P_c$ has nonzero leading
coefficient $c(c+v^2)^2$ as a polynomial in $s$. Their character
products, or equivalently their norms, are nonzero polynomials over
$\Q$. Apply Theorem~\ref{thm:density} to their finite product.
\end{proof}

The block has two free rational parameters and one monic quadratic
generator. A finite-algebra construction therefore multiplies the
existing finite rank by two and retains the same single encoder
coordinate. No additional coordinate is needed for the displayed
guards. If a source already has two parameters per pole block, this
replacement does not change its variable count.

\section{Application to prepared detector inputs}

For a detector input $b$ one may put $d=7b$. If
\[
 0<b<1/7,\quad b\text{ has square reduced denominator},\quad
 b\in\Q_2^{\times2},\quad v_2(b)\ge4,\quad v_7(b)\ge0,
\]
then $d$ satisfies \eqref{eq:prepared} and still has square reduced
denominator. The local implication for $d$ excludes negative odd
valuations of $b$ outside $\{7\}$, including at two. For the
whole-character avoidance of Proposition~\ref{prop:guards}, it is
enough that $b$ is a unit in the existing point algebra. The prepared
upper bound $b<1/7$ is only needed at the selected rational character.

These are conditions on a source's preparation. The intrinsic theorem
does not assert that a particular source has met them. A proposed
source replacement must check its complete ordered preparation and
every new guard before using the theorem.

For comparison with the scaled expression
\[
 P^{\rm sc}_{b_0}(u,s)=343b_0(7b_0+u^2)^2s^4
             -7b_0(7b_0-u^2)^2(b_0-7)^2,
\]
the exact substitutions are
\[
 b_0=7d,\qquad u=7v,\qquad
 P^{\rm sc}_{7d}(7v,s)=7^8P_d(v,s),\qquad Q_0=2401Y.
\]
Thus the small radicand is obtained by invertible rational scaling.
The accompanying exact checker verifies both conic identities, this
scaling, and the expanded coefficient table.

\begin{thebibliography}{9}
\bibitem{Burungale}
A. Burungale, \emph{A Diophantine definition of $\Z$ in $\Q$},
arXiv:2611.16130v1, 12 November 2026, Theorem B and Section 3.
\url{https://arxiv.org/abs/2611.16130}.
\end{thebibliography}
\end{document}
